\section*{Chapitre \ref{ch:b2:heat-conduction} --- Conduction thermique}

\begin{solution}{exo:b2:heat-conduction:1}
$j = \lambda\Delta T/e = \qty{80}{W/m^2}$~; \qty{8}{kW}~; avec la laine $R = 0.25 + 2.5 =
\qty{2.75}{m^2.K/W}$~: \qty{7.3}{W/m^2}, \qty{730}{W}.
\end{solution}

\begin{solution}{exo:b2:heat-conduction:2}
$a$~: \qty{1.2e-4}{m^2/s}, \qty{7e-7}{m^2/s}, \qty{1e-7}{m^2/s}. \qty{1}{cm}~:
\qty{0.9}{s}, \qty{140}{s}, \qty{1000}{s}~; \qty{20}{cm}~: \qty{6}{min}, \qty{15}{h},
\qty{5}{jours}~; croûte ($a \sim \qty{1e-6}{m^2/s}$)~: $10^{15}\,\unit{s}$, trente
millions d'années.
\end{solution}

\begin{solution}{exo:b2:heat-conduction:3}
Simple~: $R = 0.169$, $U = \qty{5.9}{W/m^2/K}$, \qty{118}{W/m^2}~; double~: $R =
0.63$, $U = \qty{1.6}{W/m^2/K}$, \qty{32}{W/m^2}. La lame d'air convecte et les vitres
échangent par rayonnement (environ la moitié du transfert réel)~; l'argon conduit
moins et convecte moins, une couche à basse émissivité supprime le
rayonnement.
\end{solution}

\begin{solution}{exo:b2:heat-conduction:4}
$L^2/a$ avec la demi-épaisseur~: \qty{27}{min}~; rôti~: $16\times$, \qty{7}{h}~;
$t \propto L^2 \propto m^{2/3}$~: $\times 1.6$.
\end{solution}

\begin{solution}{exo:b2:heat-conduction:5}
(a) $(rT')'/r = -p/\lambda$ avec $T'(0) = 0$. (b) $p = j^2/\gamma = \qty{1.7e6}{W/m^3}$~;
$pR^2/4\lambda = \qty{1e-3}{K}$. (c) \qty{625}{K}. (d) Un fil est isotherme~; le centre
d'une pastille de combustible est à des centaines de kelvins au-dessus de sa surface --- la densité
de puissance d'un réacteur est limitée par la conduction dans le combustible.
\end{solution}

\begin{solution}{exo:b2:heat-conduction:6}
(a) \cref{prop:b2:heat-conduction:wave}. (b) \qty{12}{cm}~; \qty{2.2}{m}. (c)
$\delta\ln 10 = \qty{5.2}{m}$~; retard $2.3\,\unit{rad}$, $130$ jours. (d)
L'amplitude de \qty{20}{K} doit être divisée par deux~: $x = \delta\ln 2 = \qty{1.5}{m}$.
\end{solution}

\begin{solution}{exo:b2:heat-conduction:7}
(a) $\mathrm{Bi} = hR/3\lambda \approx 10^{-4}$~; $\tau = \rho cR/3h = \qty{1100}{s}$.
(b) $\mathrm{Bi} = 1.2$~: aucune des deux limites~; le temps interne $R^2/a \approx
\qty{6000}{s}$ gouverne. (c) $R$ petit~: $\tau$ minuscule et petit \omterm{def:b2:heat-conduction:biot}{nombre de Biot}.
(d) La résistance de film $1/h$, la dominante, chute.
\end{solution}

\begin{solution}{exo:b2:heat-conduction:8}
(a) \cref{prop:b2:heat-conduction:fin}. (b) $m = \sqrt{2h/\lambda t} =
\qty{13}{m^{-1}}$, $mL = 0.39$, $\eta = 0.95$. (c) Aire $\times 2L/t = 60$, flux
$\times 57$. (d) Au-delà de $mL \approx 2$ la longueur supplémentaire est à la température
du fluide~; une ailette plus mince a un $m$ plus grand et une efficacité plus basse.
\end{solution}

\begin{solution}{exo:b2:heat-conduction:9}
(a) \cref{prop:b2:heat-conduction:resistance}. (b) Laine $\ln 2/2\pi\lambda =
\qty{2.76}{K.m/W}$, film extérieur $1/2\pi r_2h = 0.16$~: \qty{45}{W/m}~; nu,
$2\pi r_1h\Delta T = \qty{410}{W/m}$. (c) $\dd/\dd r[\ln(r/r_1)/2\pi\lambda + 1/2\pi hr]
= 0$ en $r = \lambda/h$. (d) $r_{\text{c}} = \qty{2}{cm} > \qty{1}{mm}$~: la gaine
augmente la perte --- elle refroidit le cuivre.
\end{solution}

\begin{solution}{exo:b2:heat-conduction:10}
(a) $8000(1/273 - 1/293) = \qty{2.0}{W/K}$. (b) $\sigma_s = \lambda T'^2/T^2$ avec
$T' = \qty{100}{K/m}$~; $\int\sigma_s\,\dd x = \lambda T'(1/T_{\text{c}} - 1/T_{\text{h}})
= j(1/T_{\text{c}} - 1/T_{\text{h}})$~; fois \qty{100}{m^2}~: le même. (c)
\qty{0.18}{W/K}. (d) $\Phi(1 - T_{\text{c}}/T_{\text{h}}) = \qty{550}{W}$~; rien
--- la chaleur se tient désormais à $T_{\text{c}}$ ($T_{\text{c}}\dot S_{\text{c}} = \qty{550}{W}$).
\end{solution}

\begin{solution}{exo:b2:heat-conduction:11}
(a) $j_Q(0,t) = -\lambda\partial_xT = \lambda\theta_0\sqrt2/\delta\,\cos(\omega t +
\pi/4) = \sqrt{\lambda\rho c\omega}\,\theta_0\cos(\omega t + \pi/4)$. (b) Le flux de
chaque côté au contact vaut $b_i|T_i - T_{\text{c}}|/\sqrt{\pi t}$~; en égalant~:
$T_{\text{c}} = (b_1T_1 + b_2T_2)/(b_1 + b_2)$. (c) Cuivre \qty{20.3}{\celsius},
bois \qty{29}{\celsius}. (d) Carrelage $b \approx 1500$, tapis $\approx 100$~; la
formule vaut tant que les deux corps paraissent semi-infinis~; une fois que les fronts
de chaleur atteignent l'irrigation sanguine et la face lointaine de l'objet, le régime
permanent dépend de la chaleur réellement fournie.
\end{solution}

\begin{solution}{exo:b2:heat-conduction:12}
(a) Substituer. $\theta(\pm L) = 0$ donne $\cos k_nL = 0$. (b) $a_n =
4\theta_0(-1)^n/(2n+1)\pi$~; $\theta = \sum a_n\cos(k_nx)\eu^{-ak_n^2t}$. (c) $n = 0$,
$\tau = 4L^2/\pi^2a$. (d) $\tau = \qty{650}{s}$~; centre $(4/\pi)\eu^{-t/\tau} = 0.01$~:
$t = 4.8\tau \approx \qty{52}{min}$.
\end{solution}

\begin{solution}{pb:b2:heat-conduction:1}
\textbf{1.} \qty{0.25}{m^2.K/W}~; \qty{80}{W/m^2}.

\textbf{2.} $R = 0.415$, $U = \qty{2.4}{W/m^2/K}$~; isolé $R = 2.9$, $U =
\qty{0.34}{W/m^2/K}$.

\textbf{3.} \qty{5.9}{W/m^2/K}~; \qty{1.6}{W/m^2/K}.

\textbf{4.} $4800 + 2360 = \qty{7.2}{kW}$~; $680 + 640 = \qty{1.3}{kW}$.

\textbf{5.} $UA = \qty{358}{W/K}$~: $358 \times 2500 \times 86400 = \qty{7.7e10}{J} =
\qty{21500}{kWh}$~; après, \qty{66}{W/K}~: \qty{4000}{kWh}.

\textbf{6.} $20 - 48/8 = \qty{14}{\celsius}$~; $20 - 6.8/8 = \qty{19.2}{\celsius}$~;
un mur froid rayonne moins vers le corps et peut atteindre le point de rosée.

\textbf{7.} $a = \qty{5.3e-7}{m^2/s}$~; $e^2/a = \qty{7.5e4}{s} \approx \qty{21}{h}$~: le
mur moyenne la journée (la valeur journalière $\delta = \qty{12}{cm}$ est inférieure
à son épaisseur).

\textbf{8.} Chercher $\theta_0\eu^{\iu(\omega t - \underline kx)}$~: $\underline k = (1 - \iu)/\delta$.

\textbf{9.} \qty{12}{cm}~; \qty{2.2}{m}.

\textbf{10.} $4.6\delta = \qty{54}{cm}$~; \qty{3.7}{K}~; \qty{1.1}{K}.

\textbf{11.} Un radian, $58$ jours~: mi-mars (le minimum de surface
étant à la mi-janvier).

\textbf{12.} Amplitude $\lambda\theta_0\sqrt2/\delta = \qty{6.3}{W/m^2}$, en avance de
$\pi/4$ ($45$ jours)~; cent fois le flux géothermique.

\textbf{13.} $\lambda \times \qty{0.03}{K/m} \approx \qty{0.03}{W/m^2}$, le bon ordre~; cela
fait un écart permanent de \qty{0.1}{K} sur la hauteur de la cave, noyé dans
l'oscillation saisonnière.

\textbf{14.} \qty{1e6}{W/m^2}~; $R = \qty{0.033}{K/W}$.

\textbf{15.} \qty{0.01}{K/W}~; air~: \qty{1.9}{K/W}, \qty{190}{K}.

\textbf{16.} $m = \qty{20}{m^{-1}}$, $mL = 0.61$, $\eta = 0.89$~; $\mathrm{Bi} = ht/2\lambda
= 10^{-4}$.

\textbf{17.} $A = \qty{0.108}{m^2}$~; $R_{\text{ail}} = 1/\eta hA = \qty{0.21}{K/W}$~; embase
\qty{0.006}{K/W}~; total \qty{0.26}{K/W}~: \qty{26}{K}, jonction à \qty{51}{\celsius}.

\textbf{18.} $1/hS = \qty{5.6}{K/W}$~: \qty{560}{K} --- impossible.

\textbf{19.} $R_{\text{ail}} \approx \qty{0.94}{K/W}$, total $\approx \qty{1}{K/W}$~:
\qty{125}{\celsius} --- il s'étrangle, puis s'éteint.

\textbf{20.} $V = \qty{7.2e-5}{m^3}$, \qty{0.19}{kg}, $C = \qty{175}{J/K}$~; $\tau = RC
\approx \qty{46}{s}$~; pastille~: $C = \qty{0.07}{J/K}$, $\tau \approx \qty{3}{ms}$.

\textbf{21.} La chaleur doit s'étaler de \qty{1}{cm^2} à \qty{36}{cm^2}~: le cuivre
divise par deux cette résistance d'étalement~; un caloduc déplace la chaleur par
évaporation et condensation avec une conductivité effective cent
fois celle du cuivre.

\textbf{22.} $7200(1/273 - 1/293) = \qty{1.8}{W/K}$~; \qty{0.33}{W/K}.

\textbf{23.} $100(1/298 - 1/324) = \qty{0.027}{W/K}$~; $100/324 = \qty{0.31}{W/K}$ ---
dix fois plus~: convertir du travail en chaleur est la grande irréversibilité.

\textbf{24.} $\int\lambda T'^2/T^2\,\dd x = \lambda T'\,[-1/T] = j(1/T_{\text{c}} -
1/T_{\text{h}})$.

\textbf{25.} Mur~: résistances en série, le plus mauvais conducteur commande.
Ailette~: $m = \sqrt{hp/\lambda S}$, $mL \approx 1$. Cave~: $\delta = \sqrt{2a/\omega}$,
amortissement et retard ensemble. Puce~: une chaîne de résistances de la pastille à
l'air, l'aire des ailettes faisant le travail. Partout~: la chaleur descend, et
crée de l'entropie ce faisant.
\end{solution}
